\documentclass[11pt,letterpaper]{article}

\usepackage{graphicx}           % Graphics package
\usepackage{amsmath} 
\usepackage{mathrsfs}
\usepackage[colorlinks=true]{hyperref} 
\usepackage{color}
\usepackage{amsfonts}
\usepackage[textheight=9in, textwidth=7.5in, letterpaper]{geometry}
\usepackage[makeroom]{cancel}
\usepackage{cite}
\usepackage{sectsty}
\usepackage{empheq}
\usepackage{appendix}
\usepackage{enumerate} 
\usepackage{simpler-wick}
\usepackage{amssymb}

\sectionfont{\fontsize{12}{12}\selectfont}
\subsectionfont{\fontsize{12}{12}\selectfont}

\newcommand*\widefbox[1]{\fbox{\hspace{2em}#1\hspace{2em}}}
\newcommand{\LP}{\left(}
\newcommand{\RP}{\right)}
\newcommand{\mc}[1]{\mathcal{#1}}
\newcommand{\R}{\mathbb{R}}
\newcommand{\C}{\mathbb{C}}
\newcommand{\Z}{\mathbb{Z}}
\newcommand{\D}{\partial}
\newcommand{\KET}[1]{\left| #1 \right\rangle }
\newcommand{\BRA}[1]{\left\langle #1 \right| }
\newcommand{\IP}[2]{\left\langle #1 \middle| #2 \right\rangle}
\newcommand{\PA}[2]{\frac{\partial #1}{\partial #2}}



%%%%%%%%%%%%%%%%---------------------MOST USEFUL COMMAND EVER---------------------%%%%%%%%%%%%%%%%%%%%%%
\newcommand{\fixme}[1]{{\bf {\color{red}[#1]}}}
\newcommand{\draftmode}{\usepackage[notref,notcite]{showkeys}}
\providecommand*\showkeyslabelformat[1]{\normalfont\sffamily\footnotesize#1}

%%%%%%%%%%%%%%%


\setcounter{tocdepth}{4}
\setcounter{secnumdepth}{4}

\begin{document}

\title{Physics and Category Theory}

\author{Mathew Calkins}

\date{\today}

 
\maketitle

\abstract{Extremely rough notes on physics and category theory. Maybe these will turn into something coherent, maybe not!}

\tableofcontents




\section{Fusion Categories}

\subsection{Basic Definitions}
Here we mostly follow \cite{Santaniello:2026urh}. A fusion category consists of a collection $\mc{C}$ (we will use the same symbol to refer to the category itself) plus a pair of binary operations \begin{equation}
\oplus, \otimes: \mc{C} \times \mc{C} \to \mc{C}
\end{equation}
which act roughly like the corresponding operations on familiar algebraic structures. Apparently there is a nice careful exposition of all this in \cite{Fuchs:2002cm}. For now, we just note that we have equalities\footnote{As we'll see in our favorite example, for this to be literal equality we may need the objects in $\mc{C}$ to be isomorphism classes of our favorite objects.} \begin{equation}
x \oplus y = y \oplus x, \quad x \otimes y = y \otimes x, \quad (x \oplus y) \oplus z = x \oplus (y \oplus z).
\end{equation}
So both operations are abelian, and the direct sum is associative. The product is not generally associative, but it is up to a certain natural isomorphism which we'll describe shortly. Finally, there exists a collection $\mathscr{I}$ of so-called simple objects, with the property that any object $x \in \mc{C}$ can be decomposed uniquely as a direct sum \begin{equation}
x = \bigoplus_{a \in \mathscr{I}} N^a a
\end{equation}
where necessarily\footnote{A colleague tells me that this subfield feels strongly that $0 \notin \mathbb{N}$.} $N^a \in \mathbb{Z}_{\geq 0}$. Like a basis for a vector space, this collection of simple objects is generally non-unique, though apparently its cardinality is an invariant; we call the latter the \textit{rank} of $\mc{C}$. After making of a choice of simple objects, the action of the direct product is charactized exactly by the \textit{fusion coefficients} appearing in the decompositions \begin{equation}
a \otimes b = \bigoplus_c N^{ab}_c c
\end{equation}
ranging over all triples $(a,b,c)$ of simple objects. Among the simple objects there exists a distiguished unit object $1 \in \mathscr{I}$ with the property that \begin{equation}
a \otimes 1 = 1 \otimes a = 1 \ \forall a \in \mathscr{I}.
\end{equation}
Morever, every $a \in \mathscr{I}$ has a unique dual $\bar{a} \in \mathscr{I}$ obeying \begin{equation}
N_{a \bar{a}}^1 = N_{\bar{a} a}^1 = 1.
\end{equation}

\subsection{Example: Representations of a Fixed Group}
The simplest nontrivial example of all of this arises by fixing a group $G$ and letting $\mc{C}$ consist of the representations\footnote{Modulo conjugation equivalence, if we are being careful about equality vs isomorphism} $\rho: G \to \mbox{End}(V)$ for $V$ a $\C$-linear space. Then $\oplus, \otimes$ are the usual sum and product of representations, the natural simple objects are the irreducible representations, and the fusion coefficients are Clebsch-Gordon coefficients. The unit simple object is the trivial constant scalar map \begin{equation}
g \in G \mapsto 1 \in \C^* \simeq \mbox{End}(\C)
\end{equation}
The nature of dual objects requires slightly more work. Given any irrep $\rho: G \to \mbox{End}(V)$, or in the present notation $a = (V, \rho)$, the dual object is $\bar{a} = (V^*, \rho^\vee)$, where $V^*$ is the usual dual space \begin{equation}
V^* = \{\mbox{linear maps } V \to \C\} = \mbox{End}(V,\C)
\end{equation}
and $\rho^\vee(g)$ is the inverse transpose, or equivalently acts on dual vectors by pulling back under the inverse, or equivalently \begin{equation}
\rho^\vee(g) (\lambda) = \lambda \circ \rho(g^{-1}) \in V^* \mbox{ for } \lambda \in V^*.
\end{equation}
In matrix notation, \begin{equation}
\rho^\vee(g) = (\rho(g^{-1}))^T.
\end{equation}
Now consider the product of these simple objects \begin{equation}
a \otimes \bar{a} = (V \otimes V^*, \rho \otimes \rho^\vee).
\end{equation}
The canonical evaluation map \begin{equation}
E: V \otimes V^* \to \C, \quad (v, \lambda) \mapsto \lambda(v)
\end{equation}
is $G$-invariant in this representation, in the sense that\footnote{Here we calculate only on product pairs, but by linearity this extends to general vectors in the product representation} \begin{equation}
v \otimes \lambda \xmapsto{(\rho \otimes \rho^\vee)(g)} \rho(g)(v) \otimes \rho^T(g^{-1}) (\lambda) \xmapsto{E} \LP \lambda \circ \rho(g^{-1}) \RP (\rho(g) v) = \lambda(v).
\end{equation}
More briefly, \begin{equation}
E \circ (\rho \otimes \rho^\vee)(g) = E.
\end{equation}
As already noted, finding the fusion coefficients here amounts to finding the Clebsch-Gordon coefficients. 

\bibliography{bibfile}
\bibliographystyle{utphys}

\end{document}


